\chapter{Robust Control of Sampled-Data Systems via DLMI}

\section{Introduction to Sampled-Data Systems}
Physical systems (such as aircraft, chemical reactors, and electrical grids) are continuous-time dynamical processes governed by differential equations. However, modern controllers are implemented on digital computers, microcontrollers, or distributed networks. The coupling of a continuous plant with a discrete digital controller results in a **sampled-data system**, which is hybrid in nature.

Consider the continuous-time LTI plant dynamics:
\begin{equation}
    \dot{x}(t) = A x(t) + B u(t) \label{eq:sd_plant}
\end{equation}
A digital controller samples the state vector $x(t)$ at discrete time instants $t_k$, computes the control input vector $u_k = -K x(t_k)$ using a state feedback gain $K$, and holds this control signal constant between sampling instants using a **Zero-Order Hold (ZOH)**:
\begin{equation}
    u(t) = u_k = -K x(t_k) \quad \text{for } t \in [t_k, t_{k+1}) \label{eq:sd_zoh}
\end{equation}
Substituting \eqref{eq:sd_zoh} into \eqref{eq:sd_plant} yields the hybrid closed-loop system:
\begin{equation}
    \dot{x}(t) = A x(t) - B K x(t_k) \quad \text{for } t \in [t_k, t_{k+1}) \label{eq:sd_cl}
\end{equation}
The sampling interval is denoted by $h_k = t_{k+1} - t_k > 0$. In general, the sampling intervals can be time-varying due to digital network delays, packet drops, or OS scheduling jitter:
\begin{equation}
    h_k \in [h_{\min}, h_{\max}]
\end{equation}
The presence of the term $x(t_k)$ makes the closed-loop system time-varying and infinite-dimensional, rendering standard continuous Lyapunov stability analysis ($A_{cl}^\trp P + P A_{cl} \prec 0$) inapplicable.

\section{Lyapunov Stability Analysis via Looped Functionals}
To analyze stability under variable sampling intervals, a standard quadratic Lyapunov function candidate $V(x(t)) = x(t)^\trp P x(t)$ with a constant $P \succ 0$ is highly conservative. It would require the energy to decrease monotonically at *every* continuous time instant $t$, which is unnecessary. We only require the energy to decrease from one sampling instant to the next, i.e., $V(x(t_{k+1})) < V(x(t_k))$.

To achieve this, we employ the **Looped Lyapunov Functional** approach.

Let us define a functional of the form:
\begin{equation}
    V(t) = x(t)^\trp P x(t) + \mathcal{V}_0(t), \quad t \in [t_k, t_{k+1}) \label{eq:sd_looped_lyap}
\end{equation}
where $\mathcal{V}_0(t)$ is a time-dependent looped term satisfying the boundary conditions:
\begin{equation}
    \mathcal{V}_0(t_k) = \mathcal{V}_0(t_{k+1}) = 0 \label{eq:sd_boundary}
\end{equation}
Because of the boundary conditions \eqref{eq:sd_boundary}, evaluating \eqref{eq:sd_looped_lyap} at the sampling instants yields:
\begin{equation}
    V(t_k) = x(t_k)^\trp P x(t_k) \quad \text{and} \quad V(t_{k+1}) = x(t_{k+1})^\trp P x(t_{k+1})
\end{equation}
If we can guarantee that the derivative of the functional is strictly negative throughout the sampling interval:
\begin{equation}
    \dot{V}(t) = \frac{d}{dt} \left( x(t)^\trp P x(t) + \mathcal{V}_0(t) \right) < 0 \quad \text{for all } t \in [t_k, t_{k+1})
\end{equation}
then integrating this derivative from $t_k$ to $t_{k+1}$ guarantees:
\begin{equation}
    V(t_{k+1}) - V(t_k) = x(t_{k+1})^\trp P x(t_{k+1}) - x(t_k)^\trp P x(t_k) < 0
\end{equation}
This establishes asymptotic stability at the sampling instants, which implies continuous-time stability because the states cannot escape (blow up) within a finite sampling interval.

\section{Differential LMI (DLMI) Formulation}
A powerful way to represent the looped Lyapunov functional is to allow the Lyapunov matrix to be time-varying within each sampling interval. Let $\tau(t) = t - t_k \in [0, h_k]$ represent the elapsed time since the last sample. We propose a time-dependent Lyapunov candidate:
\begin{equation}
    V(x(t), \tau(t)) = x(t)^\trp P(\tau) x(t)
\end{equation}
where $P(\tau) = P(\tau)^\trp \succ 0$ is a continuously differentiable matrix function satisfying the boundary conditions:
\begin{equation}
    P(0) = P(h_k) = P_0 \succ 0
\end{equation}
Differentiating $V(x(t), \tau(t))$ along the trajectories of \eqref{eq:sd_cl} gives:
\begin{equation}
    \dot{V}(t) = \dot{x}(t)^\trp P(\tau) x(t) + x(t)^\trp P(\tau) \dot{x}(t) + x(t)^\trp \dot{P}(\tau) x(t)
\end{equation}
where $\dot{P}(\tau) = \frac{dP(\tau)}{d\tau}$.

Substituting the dynamics $\dot{x}(t) = A x(t) - B K x(t_k)$ and defining the augmented state vector $\eta(t) = \begin{bmatrix} x(t)^\trp & x(t_k)^\trp \end{bmatrix}^\trp$, we can write $\dot{V}(t)$ as:
\begin{equation}
    \dot{V}(t) = \eta(t)^\trp \begin{bmatrix} A^\trp P(\tau) + P(\tau) A + \dot{P}(\tau) & -P(\tau) B K \\ -K^\trp B^\trp P(\tau) & 0 \end{bmatrix} \eta(t)
\end{equation}
For stability, we require this quadratic form to be negative definite, which yields the **Differential Linear Matrix Inequality (DLMI)**:
\begin{equation}
    \begin{bmatrix} A^\trp P(\tau) + P(\tau) A + \dot{P}(\tau) & -P(\tau) B K \\ -K^\trp B^\trp P(\tau) & 0 \end{bmatrix} \prec 0, \quad \forall \tau \in [0, h_k] \label{eq:sd_dlmi}
\end{equation}

\section{Numerical Solution via Multi-Vertex Parameterization}
Directly solving the DLMI \eqref{eq:sd_dlmi} is challenging because it must hold for infinitely many points $\tau \in [0, h_k]$, and it is nonlinear due to the product $P(\tau)BK$.

To render it solvable, we parameterize $P(\tau)$ as an affine function of $\tau$:
\begin{equation}
    P(\tau) = \left( 1 - \frac{\tau}{h_k} \right) P_1 + \frac{\tau}{h_k} P_2 \label{eq:sd_p_param}
\end{equation}
where $P_1 \succ 0$ and $P_2 \succ 0$ are symmetric positive definite matrices.

Taking the derivative of \eqref{eq:sd_p_param} with respect to $\tau$ yields a constant matrix:
\begin{equation}
    \dot{P}(\tau) = \frac{P_2 - P_1}{h_k}
\end{equation}
Because the DLMI is affine in $\tau$, it is guaranteed to hold for all $\tau \in [0, h_k]$ if and only if it holds at the boundary vertices $\tau = 0$ and $\tau = h_k$. This reduces the infinite-dimensional DLMI to a finite set of LMIs.

To linearize the terms involving the controller gain $K$, we perform a congruence transformation:
\begin{enumerate}
    \item Define the inverse matrices $W(\tau) = P(\tau)^{-1}$.
    \item Pre- and post-multiply the DLMI \eqref{eq:sd_dlmi} by $\text{diag}(W(\tau), W(\tau))$.
    \item Define $Y = K W_0$ where $W_0 = P_0^{-1}$.
\end{enumerate}
This yields a set of standard, convex LMIs evaluated at the vertices.

\begin{mytheorem}{Robust Sampled-Data State Feedback DLMI Synthesis}
The sampled-data hybrid system \eqref{eq:sd_cl} is stable for any variable sampling interval $h_k \in [0, h]$ if there exist symmetric matrices $W_1 \succ 0$, $W_2 \succ 0$, and a matrix $Y \in \R^{m \times n}$ satisfying the LMI system:
\begin{align}
    \begin{bmatrix} W_1 A^\trp + A W_1 - \frac{W_2 - W_1}{h} & -B Y \\ -Y^\trp B^\trp & 0 \end{bmatrix} &\prec 0 \label{eq:sd_vertex1} \\
    \begin{bmatrix} W_2 A^\trp + A W_2 - \frac{W_2 - W_1}{h} & -B Y \\ -Y^\trp B^\trp & 0 \end{bmatrix} &\prec 0 \label{eq:sd_vertex2}
\end{align}
The robust digital feedback gain is reconstructed as:
\begin{equation}
    K = Y W_1^{-1} \label{eq:sd_k_reconstruct}
\end{equation}
\end{mytheorem}
By evaluating these conditions at the limits of the sampling interval, we can use standard SDP solvers to synthesize digital controllers that are guaranteed to remain stable under time-varying, uncertain sampling rates, providing a crucial bridge between continuous physics and digital control.
